% DO NOT EDIT - automatically generated from metadata.yaml

\def \codeURL{https://github.com/DerKevinRiehl/replication_ramp_metering_rl_2020_zhouetal}
\def \codeDOI{}
\def \dataURL{}
\def \dataDOI{}
\def \editorNAME{}
\def \editorORCID{}
\def \reviewerINAME{}
\def \reviewerIORCID{}
\def \reviewerIINAME{}
\def \reviewerIIORCID{}
\def \dateRECEIVED{31 May 2026}
\def \dateACCEPTED{}
\def \datePUBLISHED{}
\def \articleTITLE{[RE] Ramp Metering for a Distant Downstream Bottleneck Using Reinforcement Learning with Value Function Approximation}
\def \articleTYPE{Replication}
\def \articleDOMAIN{Transportation Science}
\def \articleBIBLIOGRAPHY{bibliography.bib}
\def \articleYEAR{2026}
\def \reviewURL{}
\def \articleABSTRACT{
The objective of this work is to reproduce the results of the reference study \cite{hou2021rampmeter} and verify the claim that the reinforcement learning algorithms Ape-X DQN, PPO and A3C show superior performance for freeway ramp meter signal control in the metrics traffic throughput, vehicle speed, number of stops per vehicle, fuel efficiency per vehicle and per-vehicle emission rates for CO_2 and NO_x than the given baseline controllers no ramp meter, fixed-time controller and ALINEA in both mixed near-congested/congested traffic situations and extremely congested traffic situations. Furthermore, we aim to confirm that PPO shows better results in all of the aforementioned assessment criteria than the other reinforcement learning algorithms, while showing a convergence trend similar to Ape-X DQN and one that is approximately 15 times faster than A3C. To achieve this, we implemented a simulation environment in Python, together with the three baseline controllers (no ramp meter, fixed-time control and ALINEA) and the three RL-based controllers (Ape-X DQN, PPO and A3C), which we subsequently used to compare our obtained results with the original claims. 
}
\def \replicationCITE{Zhou et al. (2020). "Ramp metering for a distant downstream bottleneck using reinforcement learning with value function approximation". Journal of Advanced Transportation, 2020, Article 8813467.}
\def \replicationBIB{}
\def \replicationURL{https://doi.org/10.1155/2020/8813467}
\def \replicationDOI{10.1155/2020/8813467}
\def \contactNAME{Kevin Riehl}
\def \contactEMAIL{kriehl@ethz.ch}
\def \articleKEYWORDS{ramp metering, deep reinforcement learning, model predictive control, freeway traffic control, traffic flow theory, simulation, SUMO, Python, TensorFlow, PyTorch, MPC, DRL}
\def \journalNAME{ReScience C}
\def \journalVOLUME{4}
\def \journalISSUE{1}
\def \articleNUMBER{}
\def \articleDOI{}
\def \authorsFULL{Raphaël André Maurice Benvenuti, Krishna Kanth Vuppala Narasimha, Kevin Riehl, Anastasios Kouvelas and Michail A. Makridis}
\def \authorsABBRV{R. Benvenuti et al.}
\def \authorsSHORT{Benvenuti et al.}
\title{\articleTITLE}
\date{}
\author[1\orcid{0009-0008-4493-3722}]{Raphaël André Maurice Benvenuti}
\author[1\orcid{0009-0005-8927-8862}]{Krishna Kanth Vuppala Narasimha}
\author[1,\orcid{0000-0003-4620-8379}]{Kevin Riehl}
\author[1,\orcid{0000-0003-4571-2530}]{Anastasios Kouvelas}
\author[1,\orcid{0000-0001-7462-4674}]{Michail A. Makridis}
\affil[1]{Institute for Transportation Systems and Planning, Traffic Engineering Group, ETH Zürich, 8093 Zurich}
